% Part: sets-functions-relations
% Chapter: arithmetization
% Section: checking-details
%
\documentclass[../../../include/open-logic-section]{subfiles}


\begin{document}

\olfileid{sfr}{arith}{check}
\olsection{ক্রমিত বলয় ও ক্ষেত্র}

এই অধ্যায় জুড়ে আমরা দাবি করেছি যে কয়েকটি সংজ্ঞা ``যেমন হওয়া উচিত'' তেমন আচরণ করে। এই কারিগরি পরিশিষ্টে আমরা কথাটির অর্থ পুরোপুরি খুলে বলব এবং সংজ্ঞাগুলি যে ``সঠিকভাবে'' আচরণ করে তা (কীভাবে দেখাতে হয় তার রূপরেখাসহ) দেখাব।

\olref[int]{sec}-এ আমরা $\Int$-এর উপর যোগ ও গুণ সংজ্ঞায়িত করেছি। দেখাতে চাই যে সংজ্ঞা দুটি $\Int$-কে আমাদের ``কাঙ্ক্ষিত'' গঠনটি দেয়। বিশেষ করে, সেই গঠনটি হলো একটি বিনিমেয় বলয়ের গঠন।

\begin{defn}
	একটি \emph{বিনিমেয় বলয়} হলো একটি সেট $S$, যার সঙ্গে নির্দিষ্ট উপাদান $0$ ও $1$ এবং $+$ ও $\times$ ক্রিয়া দেওয়া থাকে এবং যা নিচের আটটি সূত্র পূরণ করে:
	\begin{align*}
		\emph{সংযোগী ধর্ম}&&a + (b+ c) & = (a + b) + c \\
		&& (a \times b) \times c & = a \times (b\times c)\\
		\emph{বিনিময় ধর্ম}&&a + b &= b+ a  \\
		&&  a \times b&= b\times a\\
		\emph{অভেদী উপাদান}&&a + 0 &= a \\
		&& a \times 1 &= a\\
		\emph{যোগাত্মক বিপরীত}&&(\exists b\in S)0&=a + b\\
		\emph{বণ্টনশীলতা}&&a \times (b+ c ) &= (a \times b) + (a \times c)
	\end{align*}
	অপ্রকাশিতভাবে, এগুলির প্রত্যেকটিই $S$-এ সীমাবদ্ধ সার্বিক পরিমাণসূচকে আবদ্ধ। আরও লক্ষ করো, এখানে $0$ ও~$1$ উপাদানদুটি একই নামের স্বাভাবিক সংখ্যা না-ও হতে পারে।
\end{defn}

তাই পূর্ণসংখ্যাগুলি একটি বিনিমেয় বলয় গঠন করে কি না যাচাই করতে শুধু এই আটটি শর্ত পূরণ হয় কি না দেখতে হবে। কোনো শর্তই প্রতিষ্ঠা করা {কঠিন} নয়, তবে কাজটি কিছুটা শ্রমসাধ্য। যেমন, যোগের ক্ষেত্রে \emph{সংযোগী ধর্ম} এভাবে প্রমাণ করা যায়:

\begin{proof}
$i, j, k \in \Int$ নির্দিষ্ট করি। তাহলে এমন $a_1, b_1, a_2, b_2, a_3, b_3 \in
\Nat$ আছে যাতে $i = \equivrep{a_1, b_1}{}$, $j =
\equivrep{a_2,b_2}{}$ এবং $k = \equivrep{a_3, b_3}{}$। (পাঠযোগ্যতার জন্য আমরা ``$\equivrep{x, y}{}$'' লিখি, ``$\equivrep{x, y}{\Intequiv}$'' নয়; এই অংশ জুড়েই তা করব।) এবার:
\begin{align*}
	i + (j + k) &= \equivrep{a_1, b_1}{}+(\equivrep{a_2, b_2}{} + \equivrep{a_3, b_3}{}) \\
	&= \equivrep{a_1,  b_1}{} + \equivrep{a_2+a_3, b_2+b_3}{}\\
	&= \equivrep{a_1 + (a_2 + a_3), b_1 + (b_2 + b_3)}{}\\
	&= \equivrep{(a_1 + a_2) + a_3, (b_1 + b_2) + b_3}{}\\
	&= \equivrep{a_1 + a_2, b_1 + b_2}{} + \equivrep{a_3, b_3}{}\\
	&= (\equivrep{a_1, b_1}{} + \equivrep{a_2, b_2}{}) + \equivrep{a_3, b_3}{}\\
	&= (i+j) + k
\end{align*}
এখানে $\Nat$-এর উপর যোগের আচরণ অবাধে ব্যবহার করেছি।
\end{proof}

একইভাবে \emph{যোগাত্মক বিপরীত} এভাবে প্রমাণ করা যায়:

\begin{proof}
$i \in \Int$ নির্দিষ্ট করি, যাতে কোনো $a,b \in
\Nat$-এর জন্য $i = \equivrep{a,b}{}$। ধরি $j = \equivrep{b,a}{} \in \Int$। স্বাভাবিক সংখ্যার আচরণ ব্যবহার করে $(a+b) + 0 = 0 + (a+b)$; তাই সংজ্ঞা অনুসারে $\tuple{a+b, b+a} \sim_\Int \tuple{0,0}$, এবং সেই কারণে $\equivrep{a+b, b+a}{} = \equivrep{0, 0}{} = 0_\Int$। অতএব এখন $i + j =
\equivrep{a,b}{}+\equivrep{b,a}{}=\equivrep{a+b, b+a}{}= \equivrep{0,
0}{} = 0_\Int$।
\end{proof}

আর \emph{বণ্টনশীলতা}-র একটি প্রমাণ এখানে:

\begin{proof}
উপরের মতো $i = \equivrep{a_1, b_1}{}$, $j =
\equivrep{a_2,b_2}{}$ এবং $k = \equivrep{a_3, b_3}{}$ নির্দিষ্ট করি। এবার:
\begin{align*}
	i \times (j + k)
	&= \equivrep{a_1, b_1}{} \times (\equivrep{a_2,b_2}{} + \equivrep{a_3, b_3}{})\\
	&= \equivrep{a_1, b_1}{} \times \equivrep{a_2 + a_3,b_2+b_3}{}\\
	&= \equivrep{a_1  (a_2 + a_3) + b_1  (b_2+b_3), a_1  (b_2 + b_3) + b_1 (a_2 + a_3)}{}\\
	&= \equivrep{a_1 a_2 + a_1a_3 + b_1 b_2+b_1b_3, a_1 b_2 + a_1b_3 + a_2b_1 + a_3b_1}{}\\
%		&= \equivrep{a_1 a_2 + b_1b_2 + a_1a_3 + b_1 b_2+b_1b_3, a_1 b_2 + a_1b_3 + a_2b_1 + a_3b_1}{}\\
	&= \equivrep{a_1a_2 + b_1b_2, a_1b_2 + a_2b_1}{} + \equivrep{a_1a_3 + b_1b_3, a_1b_3 + a_3b_1}{}\\
	&= (\equivrep{a_1, b_1}{} \times \equivrep{a_2,b_2}{}) + (\equivrep{a_1, b_1}{} \times  \equivrep{a_3, b_3}{})\\
	&= (i \times j) + (i \times k)
\end{align*}
\end{proof}

বাকি পাঁচটি শর্ত প্রমাণের কাজ অনুশীলন হিসেবে রেখে দিই। সেগুলি করার পরে আমরা দেখিয়েছি যে $\Int$ একটি বিনিমেয় বলয় গঠন করে, অর্থাৎ যোগ ও গুণ (আমাদের সংজ্ঞা অনুসারে) যেমন আচরণ করা উচিত তেমনই করে।

\begin{prob}
প্রমাণ করো যে $\Int$ একটি বিনিমেয় বলয়।
\end{prob}

কিন্তু আমাদের কাজ শেষ হয়নি। $\Int$-এর উপর যোগ ও গুণ সংজ্ঞায়িত করার পাশাপাশি আমরা একটি ক্রমসম্পর্ক $\leq$ সংজ্ঞায়িত করেছি; সেটিও যেমন আচরণ করা উচিত তেমনই করে কি না যাচাই করতে হবে। আরও বিস্তারিতভাবে, দেখাতে হবে যে $\Int$ একটি \emph{ক্রমিত} বলয় গঠন করে।\footnote{\olref[sfr][rel][ord]{def:linearorder} থেকে মনে করো, পূর্ণ ক্রম হলো এমন একটি সম্পর্ক যা প্রতিবিম্ব ধর্মবিশিষ্ট, পরিযায়ী, বিপ্রতিসম এবং সংযুক্ত। ক্রমসম্পর্কের প্রসঙ্গে সংযুক্তিকে কখনও \emph{ত্রিবিভাজন} বলা হয়, কারণ যেকোনো $a$ ও $b$-এর জন্য $a \leq b \lor a
	= b \lor a \geq b$।}

\begin{defn}
একটি \emph{ক্রমিত বলয়} হলো এমন একটি বিনিমেয় বলয়, যার সঙ্গে একটি পূর্ণ ক্রমসম্পর্ক $\leq$-ও দেওয়া থাকে এবং যা নিচের শর্তগুলি পূরণ করে:
\begin{align*}
	a \leq b &\lif a + c \leq b + c\\
	(a \leq b \land 0 \leq c) &\lif a \times c \leq b \times c
\end{align*}
\end{defn}

\begin{prob}
প্রমাণ করো যে $\Int$ একটি ক্রমিত বলয়।
\end{prob}

আগের মতোই, আমাদের নির্মিত~$\Int$ যে একটি ক্রমিত বলয় তা দেখানো শ্রমসাধ্য কিন্তু নিয়মিত কাজ। কাজটি তোমার জন্য রেখে দিই।

এতে পূর্ণসংখ্যার কাজ শেষ হলো। কিন্তু এখন মূলদ সংখ্যা সম্পর্কে খুব অনুরূপ কয়েকটি বিষয় দেখাতে হবে। বিশেষ করে $+$, $\times$ ও $\leq$-এর দেওয়া সংজ্ঞা অনুসারে মূলদ সংখ্যাগুলি যে একটি ক্রমিত \emph{ক্ষেত্র} গঠন করে তা দেখাতে হবে:
\begin{defn}\ollabel{orderedfield}
একটি \emph{ক্রমিত ক্ষেত্র} হলো এমন একটি ক্রমিত বলয়, যা আরও নিচের শর্তটি পূরণ করে:
\begin{align*}
	\emph{গুণাত্মক বিপরীত}& & (\forall a \in S \setminus \{0\})(\exists b \in S) a\times b& = 1
\end{align*}
% BN-SRC-011 TARGET-CORRECTION-BEGIN
\par\smallskip\noindent\textnormal{\small\textbf{উৎস-সংশোধন (BN-SRC-011):} হিমায়িত ইংরেজি সংজ্ঞায় ক্ষেত্রের জন্য $0\neq1$ শর্তটি বলা হয়নি। সেটি না বললে এক-উপাদানের শূন্য বলয়ও প্রদর্শিত বিপরীতের শর্ত শূন্যসত্যভাবে পূরণ করে। এখানে ক্রমিত ক্ষেত্র বলতে প্রচলিত অশূন্য ক্ষেত্র বোঝানো হয়েছে, তাই অতিরিক্তভাবে $0\neq1$ ধরা হবে।} % BN-SRC-011 TARGET-CORRECTION-END
\end{defn}

$\Int$ একটি ক্রমিত বলয় গঠন করে দেখানোর পরে $\Rat$ একটি ক্রমিত ক্ষেত্র গঠন করে তা দেখানো সহজ, কিন্তু শ্রমসাধ্য।

\begin{prob}
প্রমাণ করো যে $\Rat$ একটি ক্রমিত ক্ষেত্র।
\end{prob}

পূর্ণসংখ্যা ও মূলদ সংখ্যার কাজ শেষ করে কেবল বাস্তব সংখ্যার কাজ বাকি থাকে। বিশেষ করে দেখাতে হবে যে $\Real$ একটি \emph{পূর্ণ} ক্রমিত ক্ষেত্র, অর্থাৎ পূর্ণতা ধর্মসহ একটি ক্রমিত ক্ষেত্র গঠন করে। \olref[cuts]{realcompleteness} দেখিয়েছে যে $\Real$-এর পূর্ণতা ধর্ম আছে। তবে $\Real$ যে একটি ক্রমিত ক্ষেত্র, তা যাচাই করার (ক্লান্তিকর) ধাপগুলি এখনও শেষ করতে হবে।

ওই শ্রমসাধ্য অনুশীলনে ঝাঁপিয়ে পড়ার আগে আরও কয়েকটি ``তাৎক্ষণিক'' বিষয় যাচাই করা দরকার। যেমন, নিশ্চিত হতে হবে যে যেকোনো কর্তন $\alpha$ ও $\beta$-এর জন্য সংজ্ঞা অনুযায়ী $\alpha + \beta$ সত্যিই একটি \emph{কর্তন}। এই তথ্যটির একটি প্রমাণ এখানে:

\begin{proof}
$\alpha$ ও $\beta$ উভয়েই কর্তন বলে $\alpha + \beta = \Setabs{p
+ q}{p \in \alpha \land q \in \beta}$ হলো $\Rat$-এর একটি অশূন্য প্রকৃত উপসেট। এবার কোনো $p \in \alpha$ ও $q \in \beta$-এর জন্য $x < p + q$ ধরি। তাহলে $x - p < q$, তাই $x - p \in \beta$, এবং $x = p + (x - p) \in
\alpha + \beta$। অতএব $\alpha + \beta$ হলো $\Rat$-এর একটি প্রারম্ভিক খণ্ড। সবশেষে, যেকোনো $p + q \in \alpha + \beta$-এর জন্য $\alpha$ ও $\beta$ উভয়েই কর্তন বলে এমন $p_1 \in \alpha$ ও $q_1 \in \beta$ আছে যাতে $p < p_1$ এবং $q < q_1$; তাই $p + q < p_1 + q_1 \in \alpha +
\beta$; সুতরাং $\alpha + \beta$-র কোনো গরিষ্ঠ উপাদান নেই।
\end{proof}

অনুরূপ চেষ্টা করে তুমি যাচাই করতে পারবে যে $\alpha - \beta$, $\alpha \times \beta$ এবং $\alpha \div \beta$ কর্তন (শেষটির ক্ষেত্রে $\beta$ শূন্য-কর্তন হওয়ার ক্ষেত্রটি বাদ দিয়ে)। তবে আবারও কাজটি তোমার জন্য রেখে দিই। % BN-SRC-012 TARGET-CORRECTION-BEGIN
\par\smallskip\noindent\textnormal{\small\textbf{উৎস-ব্যাখ্যা (BN-SRC-012):} হিমায়িত উৎসের আগের কর্তন অংশে বিয়োগ ও ভাগের প্রস্তাবিত সূত্রদুটি TeX মন্তব্য হিসেবে নিষ্ক্রিয়, কিন্তু এখানে সক্রিয় পাঠে $\alpha-\beta$ ও $\alpha\div\beta$ যাচাই করতে বলা হয়েছে। এই অনুশীলন সম্পূর্ণ করতে ওই মানক ক্রিয়াদুটি আলাদাভাবে সংজ্ঞায়িত করতে হবে; বাংলা অনুবাদ উৎসের সক্রিয়/নিষ্ক্রিয় অবস্থা বদলায়নি।} % BN-SRC-012 TARGET-CORRECTION-END

\begin{prob}
প্রমাণ করো যে $\Real$ একটি ক্রমিত ক্ষেত্র।
\end{prob}

কিন্তু একটি ছোট অসমাপ্ত বিষয় গুছিয়ে নেওয়া যাক। \olref[cuts]{sec}-এ আমরা প্রস্তাব করেছি যে $\sqrt{2} =
\Setabs{p \in \Rat}{p < 0 \text{ অথবা }p^2 < 2}$ নেওয়া যায়। কিন্তু এই সেটটি যে একটি \emph{কর্তন}, তা দেখাতে হবে। তথ্যটির একটি প্রমাণ এখানে:

\begin{proof}
স্পষ্টতই এটি মূলদ সংখ্যার একটি অশূন্য প্রকৃত প্রারম্ভিক খণ্ড; তাই এর কোনো গরিষ্ঠ উপাদান নেই দেখালেই যথেষ্ট। বিশেষ করে, $p$ যদি $p^2 < 2$ শর্তপূরণকারী একটি ধনাত্মক মূলদ সংখ্যা হয় এবং $q =
\frac{2p+2}{p+2}$ হয়, তবে $p < q$ ও $q^2 < 2$ উভয়ই দেখালেই যথেষ্ট। $p < q$ দেখতে শুধু লক্ষ করো:
\begin{align*}
	p^2 &< 2\\
	p^2 + 2p &< 2 + 2p\\
	p(p + 2) &< 2 + 2p\\
	p &< \tfrac{2+2p}{p+2} = q
\end{align*}
$q^2 < 2$ দেখতে শুধু লক্ষ করো:
\begin{align*}
	p^2 &< 2\\
	2p^2 + 4p + 2 &< p^2 + 4p+ 4\\
	4p^2 + 8p + 4 &< 2(p^2 + 4p + 4)\\
	(2p+2)^2 & <2(p+2)^2\\
	\tfrac{(2p+2)^2}{(p+2)^2} &< 2\\
	q^2 &< 2
\end{align*}
\end{proof}
\end{document}
